% 人工对照底本逐节点转写；独立TikZ重绘。

\begin{tikzpicture}[x=1cm,y=-1cm,every node/.style={anchor=west,align=left,inner sep=2pt,font=\fontsize{9.5}{12.3}\selectfont},edge/.style={draw=black!65,line width=.45pt}]
\node[text width=2.900cm] (n1) at (0.000,1.281) {对实害对象及实害结果持何种心理};
\node[text width=4.794cm] (n2) at (3.320,0.275) {过失（或意外事件）：打击错误};
\node[text width=3.459cm] (n3) at (3.320,1.672) {故意（包括间接故意）};
\node[text width=4.000cm] (n4) at (7.440,1.672) {行为时对实害对象的身份有无认识错误};
\node[text width=2.123cm] (n5) at (11.860,1.056) {有：对象错误};
\node[text width=4.127cm] (n6) at (11.860,2.062) {无：无任何事实认识错误\\（可称为间接故意的类型）};
\draw[edge] (n1.east)--(3.170,1.281)--(3.170,0.275)--(n2.west);
\draw[edge] (n4.east)--(11.710,1.672)--(11.710,1.056)--(n5.west);
\draw[edge] (n4.east)--(11.710,1.672)--(11.710,2.062)--(n6.west);
\draw[edge,->] (n3.east)--(7.290,1.672)--(7.290,1.672)--(n4.west);
\draw[edge] (n1.east)--(3.170,1.281)--(3.170,1.672)--(n3.west);
\end{tikzpicture}

